%%%%%%%%%%%%%%%%%%%%%%%%%%%%%%%%%%%%%%%%%
% FRI Data Science_report LaTeX Template
% Version 1.0 (28/1/2020)
%
% Jure Demšar (jure.demsar@fri.uni-lj.si)
%
% Based on MicromouseSymp article template by:
% Mathias Legrand (legrand.mathias@gmail.com)
% With extensive modifications by:
% Antonio Valente (antonio.luis.valente@gmail.com)
%
% License:
% CC BY-NC-SA 3.0 (http://creativecommons.org/licenses/by-nc-sa/3.0/)
%
%%%%%%%%%%%%%%%%%%%%%%%%%%%%%%%%%%%%%%%%%

\documentclass[fleqn,moreauthors,10pt]{ds_report}
\usepackage[english]{babel}

\graphicspath{{./}}

\JournalInfo{Machine Learning for Data Science 1, 2025-2026}
\Archive{Homework 5 Report}
\PaperTitle{Kernel Regression Methods}
\Authors{Tjaš Ajdovec}
\affiliation{\textit{ta5946@fri.uni-lj.si}}

\makeatletter
\renewcommand{\@maketitle}{%
\twocolumn[{%
\thispagestyle{empty}%
{\includegraphics[height=2cm]{logo}}
\vskip-74pt%
{\raggedleft\small\sffamily\bfseries\@JournalInfo\\\@Archive\par}%
\vskip64pt%
{\raggedright\color{color1}\fontfamily{cmss}\fontseries{bc}\fontshape{n}\fontsize{20}{25}\selectfont \@PaperTitle\par}%
\vskip10pt%
{\raggedright\color{color1}\sffamily\fontsize{12}{16}\selectfont  \@Authors\par}%
\vskip8pt%
\begingroup%
\raggedright\sffamily\small%
\footnotesize\@affiliation\par%
\endgroup%
\vskip10pt%
}]%
}
\makeatother

\begin{document}

\flushbottom
\maketitle
\thispagestyle{empty}

\section*{Introduction}

In this homework we implemented \textbf{Kernelized Ridge Regression (KRR)} and \textbf{Support Vector Regression (SVR)} with a polynomial and an RBF kernel each. KRR minimizes a regularized squared loss in the dual space. SVR minimizes an $\varepsilon$-insensitive loss via quadratic programming \cite{smola2004tutorial}. We applied both to a 1D sine dataset for visual inspection and to a multidimensional housing dataset for MSE evaluation via cross-validation.

\section*{Part 1: Sine Dataset}

The polynomial kernel of degree $M$ is:

\begin{equation}
    \kappa(x, x') = (1 + x \cdot x')^M
\end{equation}

Expanding the binomial yields an implicit polynomial feature map up to degree $M$.

The RBF kernel with bandwidth $\sigma$ is:

\begin{equation}
    \kappa(x, x') = \exp\!\left(-\frac{\|x - x'\|^2}{2\sigma^2}\right)
\end{equation}

Its value decays with squared distance into an infinite-dimensional implicit feature space.

Both methods apply the kernel trick, replacing inner products with kernel evaluations. \textbf{KRR} finds $\alpha = (K + \lambda I)^{-1} y$ where $K_{ij} = \kappa(x_i, x_j)$, and predicts $f(x) = \sum_i \alpha_i \kappa(x_i, x)$. \textbf{SVR} adds an $\varepsilon$-insensitive tube: residuals smaller than $\varepsilon$ incur no penalty. The dual QP uses $K$ as its quadratic term, with decision variables $[\alpha_1, \alpha_1^*, \alpha_2, \alpha_2^*, \ldots]$ and box constraint $C = 1/\lambda$. Predictions are $f(x) = \sum_i (\alpha_i - \alpha_i^*)\kappa(x_i, x) + b$, where only points outside the tube (support vectors) have $\alpha_i > 0$ or $\alpha_i^* > 0$.

We applied both methods to \texttt{sine.csv} (200 noisy observations of a sinusoidal signal on $[0, 20]$). Parameters were chosen manually to obtain a good fit and a sparse SVR solution. Results are in Table~\ref{tab:sine}.

\begin{table}[h!]
    \centering
    \small
    \begin{tabular}{l l c c c c}
        \toprule
        \textbf{Method} & \textbf{Kernel} & $\sigma$/\textbf{M} & $\lambda$ & $\varepsilon$ & \textbf{Train MSE} \\
        \midrule
        KRR & RBF        & $\sigma{=}1.0$ & 0.01 & --   & \textbf{0.069} \\
        KRR & Polynomial & $M{=}5$        & 1.0  & --   & 0.977 \\
        SVR & RBF        & $\sigma{=}1.0$ & 0.01 & 0.3  & \textbf{0.072} \\
        SVR & Polynomial & $M{=}5$        & 1.0  & 0.3  & 1.020 \\
        \bottomrule
    \end{tabular}
    \caption{Sine dataset results. SVR with RBF produced 85/200 support vectors and SVR with Polynomial produced 141/200.}
    \label{tab:sine}
\end{table}

RBF fits the periodic signal well (train MSE $\approx 0.07$) and most predictions land inside the $\varepsilon$-tube, leaving only 85/200 points as support vectors. A polynomial of degree 5 cannot replicate a multi-period sine regardless of parameters, so residuals are large relative to $\varepsilon$ and 141/200 points become support vectors. The high SV count for polynomial SVR is a direct consequence of kernel mismatch, not a tuning failure. Figure~\ref{fig:sine} shows all four fits.

\begin{figure}[h!]
    \centering
    \includegraphics[width=\linewidth]{fig/kernel_methods_comparison.png}
    \caption{Sine dataset fits. Top row: KRR with RBF (left) and Polynomial (right). Bottom row: SVR with RBF (left) and Polynomial (right). Dark blue points are support vectors and the red shaded band is the $\varepsilon$-tube.}
    \label{fig:sine}
\end{figure}

\section*{Part 2: Housing Dataset}

We applied both methods to \texttt{housing2r.csv} (300 samples, 5 features). Features were standardized before computing kernel distances. We measured MSE via 5-fold CV across a grid of kernel parameter values, comparing fixed $\lambda = 1$ against the best $\lambda$ from CV over $\{0.001, 0.01, 0.1, 1, 10, 100\}$. For SVR we used $\varepsilon = 2.0$ (roughly 20\% of the target standard deviation), giving sparse solutions with 75-144 support vectors out of 300. Figure~\ref{fig:housing} shows results for all four method-kernel combinations, with support vector counts on a secondary axis for SVR.

\begin{figure}[h!]
    \centering
    \includegraphics[width=\linewidth]{fig/housing_evaluation.png}
    \caption{Housing dataset: CV MSE vs kernel parameter for KRR (top) and SVR (bottom) with polynomial (left) and RBF (right) kernels. Dashed lines show $\lambda$ selected by inner CV and solid lines use fixed $\lambda = 1$. Gray triangles (secondary axis) show support vector counts for SVR.}
    \label{fig:housing}
\end{figure}

Table~\ref{tab:housing} shows the best configurations. For KRR, RBF outperforms polynomial. For SVR, the best polynomial ($M = 2$) matches RBF almost exactly. Polynomial kernels become numerically unstable for $M \geq 4$ even after standardization, so MSE grows sharply at higher degrees despite heavy regularization. RBF peaks around $\sigma \in [2, 10]$, near the typical inter-point distance after standardization.

\begin{table}[h!]
    \centering
    \small
    \begin{tabular}{l l c c}
        \toprule
        \textbf{Method} & \textbf{Kernel} & \textbf{Best param} & \textbf{CV MSE} \\
        \midrule
        KRR & Polynomial & $M = 2$        & 27.63 \\
        KRR & RBF        & $\sigma = 10$  & \textbf{27.13} \\
        SVR & Polynomial & $M = 2$        & 25.65 \\
        SVR & RBF        & $\sigma = 2$   & \textbf{25.64} \\
        \bottomrule
    \end{tabular}
    \caption{Best CV MSE for each method and kernel on the housing dataset.}
    \label{tab:housing}
\end{table}

\section*{Discussion}

\textbf{Similarities.} Both methods use the same kernel trick and produce identical predictions in the limit $\lambda \to 0$, $\varepsilon \to 0$. Larger $\lambda$ yields a smoother, more biased fit for both.

\textbf{Differences.} SVR achieves slightly lower MSE ($\approx 25.6$ vs $\approx 27.1$ at best). The $\varepsilon$-insensitive loss discards small residuals, making SVR more robust to noise. SVR's solution is also sparse: only support vectors contribute to predictions, while KRR uses all training points. KRR solves one linear system ($O(n^3)$), while SVR requires a quadratic program and is considerably slower. SVR also adds a hyperparameter $\varepsilon$ more sensitive to mistuning than $\lambda$.

\textbf{Preference.} For this dataset, KRR is the more practical choice: the MSE gap is small ($\approx 2$ points), it is faster to fit and has one fewer hyperparameter. SVR's sparsity would matter at larger scale or when inference cost is a concern, but at $n = 300$ neither applies.

\bibliographystyle{unsrt}
\bibliography{report}

\end{document}
